\documentclass[12pt,a4paper,oneside]{ctexbook}

\usepackage[a4paper,top=30mm,bottom=25mm,left=25mm,right=25mm,headheight=15pt,headsep=5mm,footskip=15mm]{geometry}
\usepackage{setspace}
\usepackage{fancyhdr}
\usepackage{titlesec}
\usepackage{graphicx}
\usepackage{subcaption}
\usepackage{booktabs}
\usepackage{longtable}
\usepackage{array}
\usepackage{float}
\usepackage{amsmath,amssymb}
\numberwithin{equation}{chapter}
\usepackage{pgfplots}
\usepackage{tikz}
\usetikzlibrary{positioning, fit, backgrounds, arrows.meta, calc, shapes.geometric}
\usepackage{caption}
\usepackage{hyperref}
\usepackage{url}
\usepackage[superscript,compress]{cite}
\usepackage{enumitem}
\usepackage{indentfirst}
\usepackage{ulem}
\usepackage{etoolbox} % 用于一些细节调整
% 定义一个用于封面个人信息行的宏，可以自动生成对齐的下划线和标签
% #1 是标签文字，#2 是内容文字
\newcommand{\coverline}[2]{%
  \makebox[8em][l]{#1\hfill\hspace*{0.5em}} % 标签文字，固定宽度 8em（包含冒号和间距）
  \hspace*{1.0em}%
  \uline{\makebox[28em][c]{#2}}\\[1.2ex] % 内容下划线，固定宽度 28em
}

\pgfplotsset{compat=1.18}

\setlength{\parindent}{2em}
\setlength{\parskip}{0pt}
\setstretch{1.25}

\hypersetup{
  colorlinks=true,
  linkcolor=black,
  citecolor=black,
  urlcolor=blue
}

\titleformat{\chapter}{\centering\zihao{-2}\bfseries}{第\thechapter 章}{1em}{}
\titleformat{\section}{\zihao{4}\bfseries}{\thesection}{0.8em}{}
\titleformat{\subsection}{\zihao{-4}\bfseries}{\thesubsection}{0.8em}{}
\titlespacing*{\chapter}{0pt}{0.5\baselineskip}{0.5\baselineskip}
\titlespacing*{\section}{0pt}{0.5\baselineskip}{0.5\baselineskip}
\titlespacing*{\subsection}{0pt}{0.5\baselineskip}{0.5\baselineskip}

\DeclareCaptionFont{captionfont}{\songti\bfseries\zihao{-4}}
\captionsetup[figure]{font=captionfont,labelsep=space}
\captionsetup[table]{font=captionfont,labelsep=space}

\fancypagestyle{frontmatterstyle}{
  \fancyhf{}
  \fancyfoot[C]{-\roman{page}-}
  \renewcommand{\headrulewidth}{0pt}
  \renewcommand{\footrulewidth}{0pt}
}

\fancypagestyle{mainmatterstyle}{
  \fancyhf{}
  \fancyhead[C]{湖南科技大学本科生毕业设计（论文）}
  \fancyfoot[C]{-\arabic{page}-}
  \renewcommand{\headrulewidth}{0pt}
  \renewcommand{\footrulewidth}{0pt}
}

\fancypagestyle{tocmatterstyle}{
  \fancyhf{}
  \fancyhead[C]{湖南科技大学本科生毕业设计（论文）}
  \fancyfoot[C]{-\roman{page}-}
  \renewcommand{\headrulewidth}{0pt}
  \renewcommand{\footrulewidth}{0pt}
}

\begin{document}

\begin{titlepage}
  \centering
  \vspace*{2.0cm}

  % 1. 校徽 Logo
  \includegraphics[width=11cm]{hkd-logo.jpg} \par
  \vspace{1.8cm}

  % 2. 标题 - 增加字间距和粗细
  {\zihao{2}\bfseries\songti 本\hspace{0.5em}科\hspace{0.5em}毕\hspace{0.5em}业\hspace{0.5em}设\hspace{0.5em}计（论文）\par}
  \vspace{2.5cm}

  % 3. 个人信息部分
  \zihao{3}\songti
  \renewcommand{\arraystretch}{1.6} % 略微增加行高，视觉上更舒展
  
  \begin{tabular}{l@{\extracolsep{0.5em}}l} % 使用左对齐，控制列间距
    % 使用相同宽度的 makebox 确保冒号垂直对齐
    \makebox[6em][s]{题\hfill 目}： & \underline{\makebox[22em][c]{基于大语言模型的智能学习辅助系统设计与实现}} \\
    \makebox[6em][s]{作\hfill 者}： & \underline{\makebox[22em][c]{姜诏宇}} \\
    \makebox[6em][s]{学\hfill 院}： & \underline{\makebox[22em][c]{计算机科学与工程学院}} \\
    \makebox[6em][s]{专\hfill 业}： & \underline{\makebox[22em][c]{软件工程}} \\
    \makebox[6em][s]{学\hfill 号}： & \underline{\makebox[22em][c]{2205050219}} \\
    \makebox[6em][s]{校内指导老师}： & \underline{\makebox[22em][c]{余庆春}} \\
    \makebox[6em][s]{校外指导老师}： & \underline{\makebox[22em][c]{余庆春}} \\
  \end{tabular} \par

  \vfill

  % 4. 日期
  {\zihao{3}\songti 二〇二六年五月 \par}
\end{titlepage}
\thispagestyle{empty}

\cleardoublepage
\pagenumbering{roman}
\pagestyle{frontmatterstyle}

\chapter*{摘要}
\addcontentsline{toc}{chapter}{摘要}
为了解决传统教学平台存在的“资料沉淀多、过程反馈慢、课程语境弱”等问题，本文设计并开发出一个基于大预言模型的智能学习辅助系统——AI Learn。该系统主要面向老师和学生两个类型的角色，其中包括课程管理、资料管理、作业上传及批阅、课程相关的聊天指导等功能。技术架构层面使用 Next.js 和 TypeScript 实现前后端一体式开发方式，在后端统一进行 AI 相关接口调用工作。

为了提升问答的相关性和可信度，系统采用课程中心化数据结构以及检索增强生成方法，在不同的课程中将资料切分成若干部分后再进行向量匹配，在对话过程中根据一定条件选择不同资料块对回答提供依据以及规范。为保持教学连贯性，系统还提供了非流式以及流式聊天方式，以及在外部模型不可用的情况下将自动退回到演示状态。系统将认证授权、会话组织、检索插入及模型调用模块化分组，使用静态检查、单元测试、构建以及端到端测试作为质量把关手段，以便于后续的维护工作。

系统已实现包括用户界面和完整功能的Web应用产品，并做了问答测试，从测试结果来看，在本批数据集及参数设置下，系统发出的所有36次请求均被正确返回，有效率是100\%，非流式平均延迟时间为6.72秒，流式第一个响应的时间是2.73秒；RAG增强问答F1值达到0.80，明显高于无RAG情况下的0.5185。这证明了课程语境限制以及教师审核对于教学应用来说是可以同时保证性能和可靠性的，也为今后的生成式AI应用于教学平台提供了一种可行的方法。

\vspace{1em}
\noindent\textbf{关键词：}智能学习辅助；课程上下文；RAG；作业初评；教育信息化

\cleardoublepage
\chapter*{Abstract}
\addcontentsline{toc}{chapter}{Abstract}
To address the problems of traditional teaching platforms, such as excessive accumulation of learning materials, slow process feedback, and weak course context awareness, this paper designs and develops an intelligent learning assistance system based on large language models, named AI Learn. The system is mainly oriented toward two types of users, teachers and students, and provides functions including course management, material management, assignment submission and grading, and course-related conversational guidance. At the architectural level, Next.js and TypeScript are adopted to implement an integrated full-stack development approach, while AI-related API calls are centrally handled on the backend.

To improve the relevance and reliability of question answering, the system adopts a course-centered data structure and a Retrieval-Augmented Generation (RAG) approach. Course materials are segmented into multiple chunks and then vectorized for similarity matching. During conversations, different material blocks are dynamically selected under specific conditions to provide contextual support and response constraints. To maintain teaching continuity, the system supports both non-streaming and streaming chat modes, and automatically falls back to a demonstration mode when external models are unavailable. The modules for authentication and authorization, session organization, retrieval insertion, and model invocation are organized in a modular manner. Static analysis, unit testing, building, and end-to-end testing are employed as quality assurance measures to facilitate future maintenance.

The system has been implemented as a fully functional web application with a complete user interface, and question-answering evaluations have been conducted. Experimental results show that, under the current dataset and parameter settings, all 36 requests were successfully returned, achieving an effectiveness rate of 100\%. The average latency of the non-streaming mode was 6.72 seconds, while the time to first response in streaming mode was 2.73 seconds. In addition, the RAG-enhanced question-answering achieved an F1 score of 0.80, which was significantly higher than the score of 0.5185 obtained without RAG. These results demonstrate that course-context constraints and teacher supervision can simultaneously ensure both performance and reliability in educational applications, and they also provide a feasible approach for the future integration of generative AI into teaching platforms.


\vspace{1em}
\noindent\textbf{Keywords:} Intelligent Learning Support; Course Context; RAG; Assignment Pre-review; Educational Informatization

\cleardoublepage
\pagestyle{tocmatterstyle}
\tableofcontents
\cleardoublepage
\pagenumbering{arabic}
\pagestyle{mainmatterstyle}

\chapter{绪论}

\section{研究背景与问题提出}
近年来，以ChatGPT为代表的生成式AI技术正在快速发展，其在自然语言处理、内容生成和智能交互方面有很强的效果。这种技术已经开始影响高校教学，改变课堂上的知识获取、学习反馈和师生交互方式\cite{ref9,ref10}。生成式AI已经能够很好地参与到课程答疑、学习辅导以及内容生成中来，使传统的教学系统逐渐从一个信息管理工具向智能辅助平台慢慢转变。

\subsection{传统教学平台现状}
目前在高校中广泛使用的教学平台，如 Moodle、Blackboard、超星学习通以及中国大学MOOC等，已经能够支持课程发布、资料管理、在线作业以及成绩统计等基本教学流程。这类平台在“课程资源数字化”和“教学流程在线化”方面发挥了重要作用，但其核心功能仍然以资源管理和流程管理为主，而对于课程语义理解与动态教学反馈支持有限。虽然平台能够保存学生登录、学习时长、作业提交等行为数据，但这些数据大多用于教学管理与结果统计，没有进一步转化为面向学习过程的智能分析并提供个性化反馈。

\subsection{生成式AI进入教育}
随着生成式AI逐渐进入教育场景，传统教学平台开始尝试引入智能问答、自动批阅以及学习助手等功能\cite{ref1,ref5}。然而，开放域的大语言模型在教学场景中仍然有着明显的局限性：模型回答可能脱离课程资料与教学目标，产生与课堂内容不一致的结果；作业评价缺乏课程标准约束，容易导致反馈不稳定；同时，开放域模型生成结果往往缺少明确的课程依据与审核链路，教师难以追踪回答来源及生成过程，从而影响教学反馈的可控性与可信度。尤其是在大规模课程环境下，如何在提高反馈效率的同时保持教学的一致性与评价的可靠性，是智慧教学研究中的关键问题。

\subsection{当前矛盾}
高校教学当前面临的核心问题并不只是“是否拥有数字化平台”，而是“平台中的课程数据能否能真正参与教学过程”。课程资料虽然已经被上传和存储，但在实际答疑与学习反馈过程中，这些内容往往没有被有效利用，导致平台模型生成的结果与课程目标之间缺乏稳定的关联。另一方面，大班授课场景下的作业批阅工作量较大，教师难以及时地对每一位学生提供持续性的学习反馈，而学生在学习过程中更需要的是针对课程内容的阶段性指导与问题定位，而不仅仅是一次性的答案生成。

\subsection{研究问题}
基于上述背景，研究重点逐渐从“是否接入大模型”转向“如何使大模型在课程场景中稳定、可控地工作”。本文围绕课程语境下的智能学习辅助问题，设计并实现一个面向教师与学生的协同教学的智能学习辅助系统。系统重点关注课程上下文约束、检索增强生成（RAG）、权限控制、教师审核闭环以及异常回退等问题，以保证模型生成结果能够与课程内容保持一致，并满足教学场景下对于可靠性、可追溯性以及工程可维护性的要求。后续章节将分别介绍系统总体设计（第4章）、关键实现方案（第5章）以及实验验证结果（第6章）。

\section{研究目标与主要工作}

\subsection{研究目标}
本文针对高校中的课程教学场景，设计并实现一个基于生成式AI的智能学习辅助系统。系统重点关注在不同课程语境下，其教学辅助过程中的稳定性、可控性以及工程可实现性，在保证课程内容一致性与教师审核参与的前提下，为学生提供课程答疑、学习辅导以及作业反馈等功能。

与开放域通用问答系统不同，本文并不追求模型在开放场景中的内容生成能力，而是强调模型生成的结果与课程资料、教学目标以及教师要求之间要有很强的关联性。系统以大语言模型为基础，并将课程内部教学资料作为课程上下文约束，使模型生成结果尽可能保持在课程语义范围内，降低模型出现偏离课程内容或生成不稳定结果的情况。

本文还同时关注在现实场景下系统运行的工程问题，包括权限控制、异常回退、统一AI调用以及教师审核闭环等内容，以提高系统在实际运行过程中的可维护性与可扩展性。

\subsection{研究范围}
本文的研究对象主要为高校课程教学中的智能学习辅助场景，重点围绕在课程资料约束下的课程答疑，作业初评与个性化建议等的功能展开。研究内容聚焦于应用层的系统设计与工程实现，不涉及大语言模型的底层训练与参数微调，也不研究开放域搜索与通用知识生成问题。

系统中的模型能力主要依赖于现有的生成式AI接口，通过课程资料检索、上下文组织以及流程控制等方式提升模型在教学场景中的适用性。本文重点研究“如何在不同课程场景下稳定使用模型”，而不是“如何训练更强的模型”。

\subsection{主要工作}
本文完成的主要工作如下：

\begin{enumerate}
    \item 设计并实现一个面向教师与学生协同的智能学习辅助系统，完成课程管理、作业管理、学习记录以及AI辅助等核心功能模块的集成。

    \item 设计课程语境下的检索增强生成机制（Retrieval-Augmented Generation, RAG），利用课程资料构建知识上下文，提高模型的回答与课程内容之间的一致性。

    \item 实现基于生成式AI的作业初评流程，并结合教师最终评阅，使AI生成结果能够参与教学反馈，但同时保留教师对于最终评价结果的控制能力。

    \item 设计统一的AI调用与异常回退机制，在模型接口不可用或生成失败时提供降级处理，以提高系统的稳定性和可靠性。

    \item 完成系统功能测试与实验验证，从课程问答、作业反馈以及系统运行等方面对系统的效果进行分析与评估。
\end{enumerate}

\section{论文结构安排}
全文结构如图\ref{fig:thesisStructure}所示。

\begin{figure}[H]
\centering

\begin{tikzpicture}[
    font=\small,
    box/.style={
        rectangle,
        rounded corners=3pt,
        draw=black!50,
        minimum width=3.2cm,
        minimum height=1cm,
        align=center,
        fill=gray!8
    },
    phase/.style={
        rectangle,
        rounded corners=4pt,
        draw=black!60,
        minimum width=8cm,
        minimum height=1cm,
        align=center,
        fill=gray!15,
        font=\bfseries
    },
    arrow/.style={
        ->,
        thick
    }
]

% ===== 第一阶段 =====
\node[phase] (p1) at (0,0) {研究基础与问题分析};

\node[box] (c1) at (-2,-1.8) {第1章\\绪论};

\node[box] (c2) at (2,-1.8) {第2章\\相关研究};

% ===== 第二阶段 =====
\node[phase] (p2) at (0,-4) {系统分析与架构设计};

\node[box] (c3) at (-2,-5.8) {第3章\\需求分析};

\node[box] (c4) at (2,-5.8) {第4章\\总体设计};

% ===== 第三阶段 =====
\node[phase] (p3) at (0,-8) {系统实现与实验验证};

\node[box] (c5) at (-4,-9.8) {第5章\\关键实现};

\node[box] (c6) at (0,-9.8) {第6章\\测试分析};

\node[box] (c7) at (4,-9.8) {第7章\\总结与展望};

% ===== 主箭头 =====
\draw[arrow] (p1) -- (p2);

\draw[arrow] (p2) -- (p3);

\end{tikzpicture}

\caption{论文整体结构安排}
\label{fig:thesisStructure}

\end{figure} 

\section{本章小结}
本章阐述了研究背景及目的、研究价值等，并对全文内容进行了概括介绍，便于后文进一步开展相关工作进行梳理以及系统设计

\chapter{相关研究与技术基础}

\section{生成式AI教育应用研究}

\subsection{智能问答与学习辅助研究}

随着生成式AI的发展，大语言模型开始逐渐应用于高校教学中的智能问答、学习辅导以及内容生成等场景。相比传统基于规则的教学系统，生成式AI能够通过自然语言交互提供更加灵活的学习支持，从而提升学习过程中的交互性与反馈效率。

当前研究中，智能问答与学习辅助是生成式AI教育应用中的主要方向之一。杨宗凯等人分析了生成式AI对教育评价与教学活动带来的影响，认为生成式AI正在推动教学模式向智能化与个性化方向发展\cite{ref9}。张金青等人提出“学情数据+智能服务”教学模式，通过学习行为数据分析辅助教学过程\cite{ref1}。Huang等人的研究表明，生成式AI工具能够在一定程度上提高学生学习效率与课堂参与度\cite{ref5}。

现有研究说明，生成式AI在学习支持与教学反馈方面具有较强潜力，但大多数研究主要关注模型能力与交互效率，对于课程内容约束以及教学过程控制考虑相对有限。

\subsection{AI作业评价与教学反馈研究}

除智能问答外，生成式AI在作业评价与学习反馈中的应用也逐渐增多。自动评分与智能反馈能够在一定程度上减轻教师重复性工作压力，并提高大规模课程中的反馈效率。

Maher 与 Kelly 提出的“与AI共同学习”观点认为，生成式AI更适合作为学习辅助工具参与教学过程，而不应完全替代教师完成评价工作\cite{ref6}。这一观点说明，在教学评价场景中，教师仍然需要保留对于最终结果的审核与控制能力。

在高校教学环境下，作业评价不仅需要关注反馈速度，还需要保证评价标准的一致性与结果可靠性。如果完全依赖模型自动生成评价结果，可能导致课程标准不统一、反馈内容不稳定等问题。因此，如何在提升效率的同时保留教师参与，已经成为当前生成式AI教学应用中的重要研究方向。

\subsection{教学场景中的治理与可控性问题}

随着生成式AI逐渐进入实际教学场景，模型生成结果的可控性与可靠性问题开始受到关注。卢宇等人指出，生成式AI在教育中的应用应重视可解释性、伦理规范以及教学治理等问题\cite{ref10}。

开放域大语言模型虽然具备较强生成能力，但其回答可能偏离课程目标或产生与教学内容不一致的结果。在教学场景中，这种问题可能进一步影响学习反馈质量以及教师对教学过程的控制能力。此外，如果模型输出缺少明确依据与审核链路，也会降低教学结果的可追踪性。

因此，现有研究开始更加关注课程上下文约束、教师审核参与以及教学过程可监督等问题，以保证生成式AI能够在教学场景中稳定运行。

\section{RAG与学习助手系统研究}

\subsection{RAG基本原理}

\subsection{基于向量检索的增强生成 (RAG) 原理}
在本系统中，RAG 的核心在于将用户问题与课程知识库在同一向量空间内进行匹配。设课程知识库为 $D = \{d_1, d_2, \dots, d_n\}$，其中每个 $d_i$ 代表一个知识切片（Chunk）。

首先，通过嵌入模型 $E(\cdot)$ 将文本转化为高维向量。给定用户查询 $q$，系统检索逻辑可形式化为寻找相关知识子集 $\mathcal{K}$：
\begin{equation}
\mathcal{K} = \text{Top-K} \left( \left\{ \text{sim}(E(q), E(d_i)) \mid d_i \in D \right\} \right)
\end{equation}

其中，$\text{sim}(\cdot)$ 采用余弦相似度（Cosine Similarity）度量：
\begin{equation}
\text{sim}(\mathbf{u}, \mathbf{v}) = \frac{\mathbf{u} \cdot \mathbf{v}}{\|\mathbf{u}\|_2 \|\mathbf{v}\|_2}
\end{equation}

最终，生成模型 $G$ 基于原始查询 $q$ 和检索到的上下文 $\mathcal{K}$ 产生回复 $R$：
\begin{equation}
R = G(q, \text{concat}(\mathcal{K}))
\end{equation}

\subsection{RAG在教育场景中的应用}

近年来，RAG开始逐渐应用于学习助手与智能教学系统中。Prasad等人开发的 PrepMind AI 系统验证了RAG在考试备考与学习辅助场景中的应用价值\cite{ref3}。国内研究也指出，基于大语言模型的学习平台应结合课程知识库、学习行为分析以及个性化推荐等技术，以提高教学场景中的模型适用性\cite{ref4,ref8}。

在课程问答场景中，问题通常具有较强的课程依赖性，例如课程术语、作业要求以及评分标准等内容往往来源于具体教学资料。RAG能够通过检索课程内容，为模型提供与当前课程相关的上下文信息，从而减少模型生成与课程目标不一致的问题。

因此，相比于完全开放域的生成方式，RAG更适合应用于课程语境下的学习辅助系统。

\subsection{RAG系统的局限性}

虽然RAG能够增强模型回答与课程资料之间的关联性，但其效果仍然受到检索质量与上下文组织方式等因素影响。

在实际系统中，文本切分方式、向量质量、检索阈值以及上下文长度都会影响最终生成结果。当检索阈值过高时，可能导致部分有效内容无法被召回；阈值过低则容易引入大量无关信息。此外，过多的检索结果也可能导致上下文长度增加，从而影响模型生成稳定性。

因此，如何在检索准确率、上下文规模以及系统运行稳定性之间进行平衡，仍然是RAG系统设计中的重要问题。

\section{本章小结}

本章从生成式AI教育应用以及RAG学习助手研究两个方面对相关研究与技术基础进行了梳理。现有研究表明，生成式AI能够提升教学反馈效率与学习支持能力，但在课程一致性、教学可控性以及结果可靠性等方面仍然存在不足。RAG技术通过引入外部知识检索，为课程语境下的智能问答与学习辅助提供了可行方案。上述研究为后续系统需求分析与总体设计提供了理论基础与技术参考。

\chapter{系统需求分析}

\section{角色需求分析}
系统服务对象为教师与学生两类主体，角色能力边界见表\ref{tab:roleNeeds}。

在教学过程中，教师与学生的角色需求是不对等的。教师主要关心的是课程组织效率以及评价的质量，而学生更加注重的是反馈的速度以及可以获得的帮助。这种不同会导致对权限的设计也不一样。如果教师的权利不够，那么就会影响整个课程管理的过程；如果学生权利过大，那么就有可能造成跨课程之间的访问和评价不公平的问题。表3.1中这些限制就是为了防止这两种相反的风险发生，而不仅仅是简单的将两者区分开来。

为防止权限规则只是一种描述，在这里我们将需求拆分为“任务链+约束链”，任务链表示角色需要做的操作，而约束链表示操作成立需要资源限制以及时间上的要求。比如对于课程聊天来说，学生需要具备两个条件：“已入课+会话绑定课程”。再如对于作业批阅而言，老师需要满足两个条件：“课程管理权限+有查证记录”。这样可以方便后续中间件开发。

\begin{table}[H]
  \centering
  \caption{角色需求与权限边界}
  \label{tab:roleNeeds}
  \begin{tabular}{p{1.5cm}p{5.8cm}p{6.5cm}}
    \toprule
    角色 & 核心任务 & 权限边界 \\
    \midrule
    教师 & 创建课程、发布作业、查看提交、审核反馈 & 仅可管理本人课程及课程成员数据 \\
    学生 & 加入课程、学习资料、提交作业、课程聊天 & 仅可访问已加入课程及本人学习记录 \\
    \bottomrule
  \end{tabular}
\end{table}

\section{功能需求分析}
系统按照教学闭环拆分六类功能模块，见表\ref{tab:functionModules}。

模块拆分以“最短教学闭环”为标准，即系统需要包含“建课$\to$学习$\to$提交$\to$反馈$\to$追踪”，如果一个模块脱离闭环单独存在，那么体验是断层的，数据也是不能复用的。例如聊天模块如果没有课时信息的支持，则其回复是不稳定的；而提交模块如果没有审核以及记录的过程，则过程性评价是不存在的。因此表\ref{tab:functionModules} 中的六个模块是以有先后顺序进行连接而不是并列存在。反过来讲，如果先上线聊天后补充相关约束条件，那么之前的历史对话就会出现“可以查看回复但是查不到回复的原因”。

优先级也是这样。认证、课程、资料、作业、聊天是高优先级，因为他们影响到主要流程是否可行。学习记录是中优先级，但是他的数据模型也要提前设计，否则后面补充采集会带来统计口径差异问题。并不是说不重视记录的价值，而是在保证主要流程可以跑通之后再去考虑次要部分能不能扩展。

\begin{table}[H]
  \centering
  \caption{功能模块需求矩阵}
  \label{tab:functionModules}
  \begin{tabular}{p{2.6cm}p{1.6cm}p{8.5cm}}
    \toprule
    模块 & 优先级 & 功能描述 \\
    \midrule
    认证与会话 & 高 & 登录注册、Cookie/JWT 鉴权、会话状态探测 \\
    课程管理 & 高 & 教师建课、学生凭课程码加入、课程访问控制 \\
    资料管理 & 高 & 课程资料增删改查、上下文统计 \\
    作业管理 & 高 & 作业发布、学生提交、教师审核 \\
    智能聊天 & 高 & 会话管理、课程上下文问答、流式返回 \\
    学习记录 & 中 & 记录学习行为、支撑分析与追踪 \\
    \bottomrule
  \end{tabular}
\end{table}

\section{非功能需求分析}

\subsection{系统可用性需求}

系统应具备良好的可用性与基本易用性，教师与学生无需复杂配置即可完成课程学习与教学管理相关操作。课程创建、资料上传、作业提交以及AI问答等功能应保持统一的交互方式，以降低系统使用门槛。

同时，系统中的AI功能不应影响基础教学流程的正常运行。当AI服务不可用或模型调用失败时，课程浏览、资料查看以及作业管理等核心功能仍应保持可用，避免因模型异常导致整个教学流程中断。

\subsection{系统性能需求}

系统应能够满足普通教学场景下的基本性能要求。在课程问答与学习辅助过程中，系统应在合理时间内完成检索与生成过程，并返回结果，以保证交互过程的连续性。

对于课程资料上传、文本切分以及向量化等耗时操作，系统应采用异步处理方式，避免阻塞正常业务流程。在多用户同时访问情况下，系统也应保持基本稳定运行，避免出现频繁的服务异常或请求失败问题。

\subsection{数据安全与权限控制需求}

系统需要保证课程数据与用户数据的基本安全性，并对不同角色进行权限隔离。教师仅能够管理本人课程相关内容，包括课程资料、作业以及学生学习记录；学生仅能够访问已加入课程中的学习资源与个人学习数据。

由于系统中的AI问答功能需要结合课程资料进行检索，因此课程知识库访问也应受到权限控制，避免出现跨课程数据访问问题。此外，系统应对用户登录状态进行校验，并对重要操作保留必要记录，以保证系统运行过程中的安全性与可管理性。

\subsection{系统可靠性与异常处理需求}

由于系统依赖外部AI模型接口，因此需要考虑模型调用失败、检索异常以及网络错误等情况。系统应具备基本异常处理能力，在生成失败或检索失败时能够进行降级处理，避免因局部模块异常导致整体流程中断。

同时，系统应对AI调用过程中的关键错误信息进行日志记录，以便后续问题定位与系统维护。课程数据、作业记录以及学习行为数据等核心信息也应进行持久化保存，防止因系统异常导致数据丢失。

\subsection{可维护性与可扩展性需求}

系统应采用模块化设计方式，降低不同功能之间的耦合程度。AI调用、课程管理、作业管理以及检索模块应保持相对独立，以便后续功能扩展与系统维护。

在AI能力方面，系统应支持模型接口替换与统一调用管理，使后续能够根据需求接入不同类型的大语言模型或检索服务。此外，系统后续还应具备扩展更多教学辅助功能的能力，例如学习行为分析、个性化推荐以及多课程知识支持等。

\subsection{教学可控性需求}

教学场景中的AI系统不仅需要提供生成能力，还需要保证教学过程中的可控性。系统中的AI生成结果不应直接替代教师完成教学评价，而应作为辅助参考参与教学流程。教师需要保留对于作业评价以及教学反馈结果的最终审核权。

同时，系统中的AI问答与学习辅助功能应尽可能基于课程资料与课程上下文生成内容，以减少回答偏离教学目标的情况。对于AI参与生成的结果，系统还应保留必要的过程记录与行为数据，以提高教学反馈过程的可追踪性与可管理性。

\begin{table}[H]
\centering
\caption{系统非功能需求指标}
\label{tab:nonfunctional}
\begin{tabular}{p{4cm} p{10.5cm}}
\toprule
非功能需求 & 系统要求 \\
\midrule

系统可用性 &
AI模块异常时不影响课程浏览、资料查看与作业管理等核心教学功能 \\

系统性能 &
课程问答与资料检索应在合理时间内完成，避免长时间阻塞用户交互 \\

数据安全与权限控制 &
不同角色用户应具备独立访问权限，课程资料需基于课程范围进行访问控制 \\

系统可靠性 &
系统应具备异常处理与日志记录能力，并支持AI调用失败时的降级处理 \\

可维护性与可扩展性 &
AI调用、课程管理与检索模块应采用模块化设计，支持后续模型与功能扩展 \\

教学可控性 &
教师应保留教学评价最终审核权，AI生成结果需尽可能基于课程上下文生成 \\

\bottomrule
\end{tabular}
\end{table}

\section{系统可行性分析}

\subsection{技术可行性}

本文系统基于现代 Web 技术栈与现有生成式AI接口实现，整体技术方案具有较好的工程可实现性。系统采用 Next.js 实现前后端一体化开发，统一完成页面渲染、路由管理以及后端业务逻辑处理；数据层使用 SQLite 与 Prisma 进行存储和对象映射，以便于后续的数据维护、功能扩展与数据库迁移。相关技术已经广泛应用于实际 Web 系统开发，具有较成熟的生态支持与开发文档。

在AI能力方面，系统主要通过调用现有大语言模型接口实现课程问答与作业辅助功能，不涉及底层模型训练与参数微调，因此能够降低系统开发复杂度。课程资料检索部分采用RAG（Retrieval-Augmented Generation）方案，通过向量检索增强模型对于课程内容的理解能力。当前向量数据库、Embedding模型以及相关检索框架均已经具备较成熟的实现方案，能够满足课程场景中的知识检索需求。

此外，系统整体采用模块化设计方式，AI调用、课程管理、检索流程以及用户管理等模块之间相对独立，便于后续系统维护与功能扩展。因此，从开发环境、技术成熟度以及系统实现角度来看，本文系统具有较好的技术可行性。

\subsection{系统运行可行性}

系统主要面向高校课程教学场景，整体运行流程与传统教学平台保持一致，不需要教师与学生改变原有教学习惯。教师仍然负责课程资料上传、作业布置以及教学评价等核心工作，AI模块主要作为辅助工具参与课程问答与作业初步反馈过程。

相比完全依赖AI自动完成教学评价的方式，本文系统保留教师审核与人工干预环节，使AI生成结果能够作为教学辅助参考，而非直接替代教师完成教学决策。这种方式能够在提高反馈效率的同时，保持教学过程中的可控性与稳定性。

同时，系统中的AI能力采用增量式接入方式。当AI服务不可用或生成失败时，课程浏览、资料查看以及作业管理等基础功能仍可正常运行，不会影响系统核心教学流程。因此，系统在高校教学环境中具有较好的运行可行性。

\subsection{数据与流程可行性}

高校教学平台本身已经积累了大量课程资料、作业内容以及学习行为数据，这些数据能够为生成式AI辅助教学提供基础知识来源。课程讲义、实验文档、作业要求以及评分标准等内容均可以作为RAG系统中的课程知识库，从而为模型生成提供课程上下文支持。

与此同时，高校教学过程本身具备较明确的流程结构，包括课程发布、资料学习、作业提交以及教师反馈等环节。本文系统并未改变原有教学流程，而是在现有流程基础上增加AI辅助问答与学习支持功能，因此能够较自然地融入已有教学环境。

此外，教师审核机制与学习行为记录也能够进一步形成教学反馈闭环，使系统不仅具备内容生成能力，同时能够保留必要的教学过程记录与学习数据，为后续教学分析与系统优化提供基础。

\subsection{AI应用可控性分析}

生成式AI在教学场景中的应用不仅需要关注生成能力，还需要考虑教学过程中的可靠性与可控性问题。开放域大语言模型虽然能够生成较自然的回答，但如果缺少课程约束与审核机制，容易产生与课程内容不一致或缺少依据的结果。

为提高系统在教学场景中的可控性，本文系统采用课程资料约束与教师审核结合的方式。课程问答过程主要基于课程知识库完成检索增强生成，使模型回答尽可能与课程内容保持一致；在作业评价场景中，AI仅负责生成初步反馈结果，教师仍保留最终审核与修改权。

此外，系统还保留AI调用日志、学习行为记录以及关键操作信息，以提高教学反馈过程的可追踪性。当模型生成失败或检索异常时，系统能够进行降级处理，避免AI模块异常影响整体教学流程。因此，系统在生成式AI应用过程中具备一定的可控性与工程稳定性。

\section{本章小结}
本文在本章总结角色、功能以及非功能性需求并提出相应的评估标准，并在三类典型场景上对需求做了进一步细化；最后从技术、经济与操作三个维度核对了方案的可行性，为之后的总体设计及实现过程提供指导。

\chapter{系统总体设计}

本章在前文需求分析基础上，对 AI Learn 智能学习辅助系统的整体结构进行设计。系统总体设计重点关注课程语境下的AI问答、作业辅助评价以及教学过程中的可控性问题，并从系统架构、功能模块、RAG问答流程以及异常处理机制等方面展开说明，为后续关键模块实现提供整体设计依据。

\section{系统总体架构设计}

AI Learn 系统整体采用前后端分离架构，并结合生成式AI接口与RAG检索机制实现课程学习辅助功能。系统主要由前端交互层、业务服务层、AI服务层以及数据存储层构成。

其中，前端负责课程展示、作业管理以及AI问答交互；业务服务层负责课程管理、权限控制与作业流程处理；AI服务层负责模型调用、课程资料检索以及提示词编排；数据层负责课程资料、学习记录与用户数据存储。

系统总体架构如图 \ref{fig:systemArchitecture} 所示。

\begin{figure}[H]
\centering
\resizebox{0.95\textwidth}{!}{
\begin{tikzpicture}[
    font=\small,
    arrow/.style={->, >=Latex, thick},
    module/.style={draw, rounded corners=4pt, align=center, minimum width=3.8cm, minimum height=1cm},
    core/.style={draw, rounded corners=4pt, align=center, minimum width=3.8cm, minimum height=1cm, fill=gray!15},
    external/.style={draw, dashed, rounded corners=4pt, align=center, minimum width=3.8cm, minimum height=1cm}
]
% --- 节点定义 (确保这些指令之间没有完全的空行) ---
\node[module] (user) {学生 / 教师用户};
\node[module, below=1cm of user] (frontend) {前端交互层\\Dashboard / AI聊天 / 作业页面};
\node[module, below=1cm of frontend] (service) {业务服务层\\课程管理 / 权限控制 / 作业管理};
\node[core, below left=1.5cm and 0.2cm of service] (ai) {AI服务层\\Prompt编排 / RAG检索};
\node[module, below right=1.5cm and 0.2cm of service] (database) {数据存储层\\PostgreSQL / 学习记录};
\node[module, below=1.2cm of ai] (vector) {向量检索模块\\课程资料切片 / 相似度检索};
\node[external, below=1.2cm of database] (llm) {外部大模型API\\OpenAI / DeepSeek};
% --- 连线 ---
\draw[arrow] (user) -- (frontend);
\draw[arrow] (frontend) -- (service);
\draw[arrow] (service) -- (ai);
\draw[arrow] (service) -- (database);
\draw[arrow] (ai) -- (database);
\draw[arrow] (ai) -- (vector);
\draw[arrow] (ai) -- (llm);
\draw[arrow] (database) -- (vector);
\end{tikzpicture}
}
\caption{AI Learn 系统总体架构}
\label{fig:systemArchitecture}
\end{figure}

系统整体以课程数据为核心，通过RAG检索增强模型对于课程内容的理解能力，同时保留教师审核与权限控制机制，以保证AI辅助过程能够符合教学场景中的稳定性与可控性要求。

\section{系统功能模块设计}

根据系统需求，AI Learn 主要划分为用户管理模块、课程管理模块、作业管理模块以及AI学习助手模块。系统功能结构如图 \ref{fig:functionModules} 所示。

\begin{figure}[H]
\centering

\begin{tikzpicture}[
box/.style={
draw,
rounded corners,
minimum width=3.2cm,
minimum height=1cm,
align=center,
font=\small
},
arrow/.style={->, thick}
]

\node[box] (core) at (0,0) {AI Learn系统};

\node[box] (user) at (-5,-2.2) {用户管理模块};

\node[box] (course) at (-1.7,-2.2) {课程管理模块};

\node[box] (assignment) at (1.7,-2.2) {作业管理模块};

\node[box] (ai) at (5,-2.2) {AI学习助手模块};

\draw[arrow] (core) -- (user);
\draw[arrow] (core) -- (course);
\draw[arrow] (core) -- (assignment);
\draw[arrow] (core) -- (ai);

\end{tikzpicture}

\caption{系统功能模块结构}
\label{fig:functionModules}

\end{figure}

各模块职责如下：

\begin{table}[H]
\centering
\caption{系统主要功能模块}
\label{tab:modules}
\begin{tabular}{p{4cm} p{9cm}}
\toprule
功能模块 & 模块职责 \\
\midrule

用户管理模块 &
负责教师与学生身份认证、权限校验以及用户信息管理 \\

课程管理模块 &
负责课程创建、课程加入、资料上传以及课程内容管理 \\

作业管理模块 &
负责作业发布、提交、批阅以及学习反馈记录 \\

AI学习助手模块 &
负责课程问答、RAG检索、Prompt编排以及AI辅助生成 \\

\bottomrule
\end{tabular}
\end{table}

其中，AI学习助手模块是系统核心模块，其主要实现流程将在后续章节中进一步展开。

\section{系统异常处理与安全设计}

由于系统依赖外部AI服务，因此需要考虑模型调用失败、检索异常以及网络错误等情况。系统在设计过程中采用异常回退与日志记录机制，以保证整体教学流程稳定运行。

当模型生成失败、向量检索结果为空或外部AI服务不可用时，系统不会直接中断当前课程流程，而是返回基础提示信息、普通课程内容或默认结果以维持系统基本功能运行。该方式属于系统降级处理机制，其主要目的是避免局部AI模块异常影响整体教学使用体验。

在权限控制方面，系统基于用户角色对课程资料与学习数据进行隔离，不同课程之间的数据访问受到限制，从而保证课程内容与学习记录的基本安全性。

\section{数据库与数据流设计}

系统主要数据包括用户信息、课程资料、作业记录、学习行为以及AI问答记录等内容。数据库采用关系型结构进行组织，并结合向量索引完成课程资料检索。

系统核心数据关系如图 \ref{fig:databaseRelation} 所示。

\begin{figure}[H]
\centering
\begin{tikzpicture}[
    font=\small,
    % 实体：矩形
    entity/.style={draw, thick, rectangle, minimum width=2.2cm, minimum height=0.7cm, fill=blue!5},
    % 联系：菱形
    relationship/.style={draw, thick, diamond, aspect=2.5, minimum width=1.8cm, fill=yellow!5, inner sep=1pt},
    % 属性：椭圆
    attribute/.style={draw, thin, ellipse, minimum width=1cm, minimum height=0.5cm, fill=white, font=\scriptsize},
    line/.style={draw, thick},
    card/.style={font=\scriptsize, inner sep=2pt}
]

% --- 第一层：用户与课程 ---
\node[entity] (user) at (0,0) {用户 (User)};
\node[attribute] (u_id) at (-1.8, 0.5) {\underline{id}};
\draw[thin] (user) -- (u_id);

\node[relationship] (member) at (3.5, 0) {选修/管理};

\node[entity] (course) at (7,0) {课程 (Course)};
\node[attribute] (c_id) at (8.8, 0.5) {\underline{id}};
\draw[thin] (course) -- (c_id);

% --- 第二层：中间联系层 ---
\node[relationship] (chat) at (0, -2.5) {发起};
\node[relationship] (belong) at (7, -2.5) {发布};

% --- 第三层：具体业务实体 ---
\node[entity] (session) at (0, -5) {对话记录 (Session)};
\node[entity] (assign) at (7, -5) {作业 (Assignment)};

% --- 核心逻辑连线 ---
% 用户 <-> 课程 (M:N)
\draw[line] (user) -- (member) node[card, pos=0.2, above] {m};
\draw[line] (member) -- (course) node[card, pos=0.8, above] {n};

% 用户 <-> 对话 (1:N)
\draw[line] (user) -- (chat) node[card, pos=0.2, left] {1};
\draw[line] (chat) -- (session) node[card, pos=0.8, left] {n};

% 课程 <-> 作业 (1:N)
\draw[line] (course) -- (belong) node[card, pos=0.2, right] {1};
\draw[line] (belong) -- (assign) node[card, pos=0.8, right] {n};

% --- 可选：如果需要展示行为记录，可以像这样简单穿插一条线 ---
\node[relationship] (record) at (3.5, -3.5) {产生记录};
\draw[line] (user) -- (record) node[card, pos=0.3, sloped, above] {1};
\draw[line] (record) -- (assign) node[card, pos=0.7, sloped, above] {n};

\end{tikzpicture}
\caption{系统核心数据实体关系 (ER) 图}
\label{fig:databaseRelation}
\end{figure}

数据库设计既需要满足课程业务管理需求，也需要支持后续AI检索与学习行为分析等功能扩展。

\section{本章小结}

本章对 AI Learn 智能学习辅助系统的总体结构进行了设计，并从系统架构、功能模块、RAG课程问答、AI辅助评价以及异常处理等方面说明系统整体组织方式。相关设计为后续关键模块实现与系统测试分析提供了整体结构基础。

\chapter{关键模块设计与实现}

\section{AI 服务调用链路设计}
在系统的后端架构中，AI 服务的调用并非简单的 API 请求，而是一个经过多层中间件处理的流水线（Pipeline）。为了保证请求的原子性与可扩展性，本文将对话生成过程抽象为状态 $S$ 的演进过程。

设初始状态 $S_{init} = \{U, Q\}$，其中 $U$ 代表用户信息上下文，$Q$ 为用户当前提问。整个调用链路可细化为以下算子操作：
\begin{enumerate}
    \item \textbf{状态增强算子 ($f_{ext}$)}：
    通过 Session ID 检索 Redis 中的历史对话记录 $H$，并从向量数据库中提取关联课程切片 $C$：
    \begin{equation}
    S_{enriched} = f_{ext}(S_{init}) = \{U, Q, H, C\}
    \end{equation}
    
    \item \textbf{提示词工程算子 ($f_{prompt}$)}：
    将增强后的状态映射为满足 LLM 输入规范的结构化 Prompt。设 $T_{sys}$ 为预设的教学助手角色模板：
    \begin{equation}
    P = f_{prompt}(S_{enriched}, T_{sys})
    \end{equation}
    
    \item \textbf{流式输出算子 ($f_{stream}$)}：
    调用大模型接口并建立 Server-Sent Events (SSE) 连接。
\end{enumerate}

综上，系统的核心逻辑可表达为复合函数 $\mathcal{R} = f_{stream} \circ f_{prompt} \circ f_{ext}(S_{init})$。这种函数式设计使得系统能够灵活地在中间环节插入敏感词过滤、Token 统计等监控逻辑。

\section{课程上下文 RAG 编排}
RAG 编排过程如图\ref{fig:ragFlow}，先检索再注入，检索不可用时退回到整个文档的内容。

\begin{figure}[H]
\centering
\resizebox{\textwidth}{!}{
\begin{tikzpicture}[
    node distance=1cm and 1.5cm,
    >=stealth,
    thick,
    base/.style={draw, rounded corners=2pt, align=center, minimum height=1.1cm, font=\small, inner sep=6pt},
    input/.style={base, fill=blue!5},
    proc/.style={base, fill=gray!5},
    core/.style={base, fill=orange!10, draw=orange!70, thick},
    arrow/.style={->, thick, color=gray!80}
]

    % --- 节点定义 ---
    \node[input] (q) {用户问题};
    \node[proc, right=of q] (e) {Embedding\\问题向量化};
    \node[core, right=of e] (v) {LanceDB\\向量检索 (Top-K)};
    
    % 这里的定位将 f 往右挪，d 往下挪，防止拥挤
    \node[proc, right=2cm of v] (f) {阈值过滤\\与内容截断};
    \node[proc, below=1.8cm of f] (d) {策略降级\\全文资料上下文};
    \node[input, right=1.5cm of f] (p) {Prompt 注入\\LLM 调用};

    % --- 连线逻辑 ---
    \draw[arrow] (q) -- (e);
    \draw[arrow] (e) -- (v);
    
    % 主路径：文字置于线上方
    \draw[arrow] (v) -- node[above, font=\scriptsize, black] {检索成功} (f);
    \draw[arrow] (f) -- (p);
    
    % 降级路径修复：
    % 1. 使用 pos=0.5 将文字放在折线的转角处之后
    % 2. 使用 anchor=north 强制文字在连线正下方，不往节点上靠
    \draw[arrow] (v) |- (d) 
        node[pos=0.6, below, font=\scriptsize, color=red!70!black, align=left] {相似度不足 / 结果为空};
    
    % 3. 降级结果指向 Prompt 注入
    \draw[arrow] (d) -| (p);

\end{tikzpicture}
}
\caption{课程上下文 RAG 编排流程}
\label{fig:ragFlow}
\end{figure}

在提示词构造阶段，系统引入上下文预算约束：
\begin{equation}
\sum_{i=1}^{n}\lvert c_i \rvert \le C_{\max},
\label{eq:contextBudget}
\end{equation}
其中 $c_i$ 为第 $i$ 个注入片段，$C_{\max}$ 为最大字符预算，防止上下文过多造成时延抖动。

RAG编排实现可以总结为五个步骤：问题预处理、向量化、候选检索、片段筛选、提示词注入。首先进行长度裁剪以及语义清洗，防止不合逻辑的问题造成向量化代价过高；然后做课程范围内的Top-K检索；接着根据相似度大小进行降序排列；最后以字符数为单位逐个插入片段，构成课程约束提示。其中最重要的一个点是“课程内检索”，即所有的候选片段都必须有相同的课程ID，防止不同课程的知识混淆。而Top-K数量、相似度阈值以及字符数等参数都在lib/chat-rag-retrieval.ts中定义，在lib/chat-pipeline.ts中的会话编排处进行提示词注入拼接。调参过程中也发现两种情况，即阈值较低则召回片段增多但是噪声增加；而阈值较高则与课程相关的一些边缘化知识被忽略。所以并不是一个最优阈值，而是在“可解释性和时延之间的权衡”  。

在参数上，RAG主要受到Top-K、最小相似度阈值、最大注入片段数量、最大注入字符数的影响。Top-K表示召回上限，阈值表示对于噪声的过滤程度，片段数以及字符数对时延和上下文密度都有影响，在保证服务可用性的前提下，系统设有两层退化策略：如果向量表不存在或者向量化失败以及检索未命中，则切换到全量文档注入；如果全量文档仍然不足，则模型会给出“资料未覆盖该问题”的提示。这个退化路径主要是为了解决“检索端失败不应该影响主线程”的工程需求，并不是希望每个问题都能得到高置信的回答。

\section{作业初评与教师审核闭环实现}
作业评审闭环是本系统的核心功能之一，旨在通过“AI 预审+教师终审”的协作模式提升教学效率。该流程的逻辑架构如图\ref{fig:reviewFlow}所示。

\begin{figure}[H]
  \centering
  \begin{tikzpicture}[node distance=8mm and 10mm, every node/.style={draw, rounded corners, align=center, minimum width=2.8cm, minimum height=1.1cm, fill=white, font=\small}]
    % 节点定义
    \node (s1) {学生提交作业 \\ \footnotesize (Payload)};
    \node[right=of s1] (s2) {AI 语义初评 \\ \footnotesize (Reasoning)};
    \node[right=of s2] (s3) {持久化存储 \\ \footnotesize (StudyRecord)};
    \node[right=of s3] (s4) {教师多维审核 \\ \footnotesize (Human-in-the-loop)};
    
    % 连线
    \draw[->, thick, >={Stealth}] (s1) -- (s2);
    \draw[->, thick, >={Stealth}] (s2) -- (s3);
    \draw[->, thick, >={Stealth}] (s3) -- (s4);
    
    % 修正后的标注说明 (使用 \itshape 代替不存在的 italic key)
    \node[below=0.4cm of s2, draw=none, fill=none, font=\scriptsize\itshape] {输出项：strengths, issues, \\ suggestions, scoreSuggestion};
  \end{tikzpicture}
  \caption{作业初评与教师审核闭环流程图}
  \label{fig:reviewFlow}
\end{figure}

\subsection{自适应评分融合模型}
为了平衡自动化评估的效率与人工审核的权威性，本文设计了一个自适应评分融合模型。最终得分 $S_{\text{final}}$ 不再由单一主体决定，而是基于教师的投入度与 AI 的置信度进行动态加权：

\begin{equation}
S_{\text{final}} = \beta(\Delta, \tau, \rho) \cdot S_{\text{ai}} + \bigl(1 - \beta(\Delta, \tau, \rho)\bigr) \cdot S_{\text{teacher}} - \lambda \cdot \mathbb{1}_{\text{rubber-stamp}}
\label{eq:scoreFusion}
\end{equation}

其中，$S_{\text{ai}}$ 与 $S_{\text{teacher}}$ 分别代表 AI 建议分与教师评分。模型引入了三个关键特征变量：分歧度 $\Delta = |S_{\text{teacher}} - S_{\text{ai}}|$、审核耗时 $\tau$ 以及 AI 信息完整度 $\rho$。

\subsubsection{权重自适应衰减算子}
权重系数 $\beta$ 并非固定常数，其随审核环境特征自适应演进。其构造逻辑遵循“教师干预越多，AI 权重越低”的负反馈原则：
\begin{equation}
\beta(\Delta, \tau, \rho) = \beta_0 \cdot \rho \cdot \exp\left( -\alpha_1 \frac{\Delta}{\Delta_0} - \alpha_2 \frac{\tau_0}{\tau + \tau_0} \right)
\label{eq:scoreFusionBeta}
\end{equation}
在此模型中，$\beta_0 \le 0.4$ 设定了 AI 影响力的物理上限，确保教学责任的最终主体始终为教师。$\rho \in [0,1]$ 为 AI 输出信息的质量权重，由 \texttt{lib/ai.ts} 输出的结构化字段长度综合计算得到。指数项通过分歧度与耗时的乘性衰减，实现了权重的非线性平滑过渡。

\subsubsection{“鸭子盖章”行为的识别与惩罚}
为了防止教师产生过度依赖 AI 的“惯性审美”或为了节省时间而盲目通过（即“鸭子盖章”行为），模型引入了示性惩罚项 $\mathbb{1}_{\text{rubber-stamp}}$：
\begin{equation}
\mathbb{1}_{\text{rubber-stamp}} = 
\begin{cases} 
1, & \text{if } \Delta=0 \land \tau < \tau_{\min} \land \rho < \rho_{\min} \\ 
0, & \text{otherwise} 
\end{cases}
\label{eq:rubberStamp}
\end{equation}
当系统检测到教师在 AI 信息完整度较低（$\rho < \rho_{\min}$）且未经过充分审阅（$\tau < \tau_{\min}$）的情况下直接接受 AI 分数时，系统将触发惩罚因子 $\lambda$，并将该记录标记为待二次抽检状态。

\subsection{闭环逻辑的分段场景分析}
基于式\eqref{eq:scoreFusion} 至式\eqref{eq:rubberStamp}，系统在实际应用中覆盖了以下典型教学场景：
\begin{enumerate}[label=(\arabic*)]
    \item \textbf{深度人工介入（教师细评）}：当教师对作业进行详尽修改（$\tau \gg \tau_0$）或大幅修正分值（$\Delta$ 较大）时，$\beta$ 随指数项快速衰减。此时 $S_{\text{final}}$ 几乎完全由教师意志主导，AI 仅作为提供初步线索的工具。
    \item \textbf{中度协同校准（小幅修改）}：在常规审核中，$\beta$ 处于中位区间，模型在两方意见间进行平滑插值，这与原始线性融合的行为一致，能有效过滤 AI 因理解偏差产生的评分波动。
    \item \textbf{异常行为防御（快速同意）}：若 AI 因学生提交内容不足而给出低置信度建议（$\rho$ 较低），此时教师若快速点击通过（$\Delta=0, \tau$ 极小），则会触发惩罚项。这不仅纠正了最终得分，还通过数据审计迫使教学过程回归严谨。
\end{enumerate}

本闭环机制的设计目标并非取代教师，而是通过结构化的预审（strengths、issues、suggestions）节约教师检索学生错误点的时间。所有评审结果最终写入 \texttt{StudyRecord} 的元数据中，确保了每一笔评分记录均可追溯、可审计。

\section{聊天时序与会话管理实现}
课程聊天时序如图\ref{fig:chatSequence}。会话管理接口支持 GET/POST/PATCH/DELETE，实现会话创建、模型切换、重命名和删除。

会话管理采用显式生命周期：创建（POST）- 使用（GET/聊天）- 更新（PATCH）- 删除（DELETE）。创建时绑定课程 ID 和模型配置，以便后续问答具有课程界限；更新时可以更改标题以及模型但是会话所有者不会发生变化；删除时级联删除相关消息从而防止出现孤立条目。这使得会话状态转移更加可靠并且便于审核及回滚等操作。

时序方面更关注幂等性和并发一致性而非把流程图画得美观。重复提交情况下，接口利用会话以及消息主键防止重复写入，在此过程中出现的问题是：前端网络抖动导致重试会导致相同的消息被再次提交，如果没有相应的限制条件，旧的消息会被重新写入并且影响到接下去的模型输入；流式断开后，在finally部分统一关闭流并返回结束标志使前端的状态机收敛；产生错误之后回滚该次用户的写入以保证消息的顺序以及可见回复的一致性，防止对接下来的推理造成干扰。

\usetikzlibrary{arrows.meta, positioning}

\begin{figure}[H]
  \centering
  \begin{tikzpicture}[
    node distance=2.5cm,
    line/.style={-Stealth, thick},
    lifeline/.style={dashed, gray},
    every node/.style={font=\small}
  ]
    % 定义角色
    \node (u) {用户前端};
    \node (a) [right=of u] {聊天接口};
    \node (p) [right=of a] {聊天流水线};
    \node (db) [right=of p] {数据库/检索层};

    % 画生命线
    \foreach \x in {u,a,p,db} {
      \draw[lifeline] (\x) -- ++(0,-5);
    }

    % 画交互消息（使用 yshift 统一偏移量）
    \draw[line] ([yshift=-0.8cm]u.center) -- node[above]{1. 发送消息} ([yshift=-0.8cm]a.center);
    \draw[line] ([yshift=-1.6cm]a.center) -- node[above]{2. 鉴权与校验} ([yshift=-1.6cm]p.center);
    \draw[line] ([yshift=-2.4cm]p.center) -- node[above]{3. 读取上下文} ([yshift=-2.4cm]db.center);
    \draw[line] ([yshift=-3.2cm]db.center) -- node[above]{4. 返回资料} ([yshift=-3.2cm]p.center);
    \draw[line] ([yshift=-4.0cm]p.center) -- node[above]{5. 生成并落库} ([yshift=-4.0cm]a.center);
    \draw[line] ([yshift=-4.8cm]a.center) -- node[above]{6. 返回响应} ([yshift=-4.8cm]u.center);

  \end{tikzpicture}
  \caption{课程聊天核心时序}
  \label{fig:chatSequence}
\end{figure}

\section{系统界面展示}

为了便于用户完成课程学习、AI问答以及作业管理等操作，系统提供了控制台主页、AI学习助手页面以及教师审核页面等交互界面。界面设计以教学流程为中心，尽量减少复杂操作，提高课程学习过程中的交互效率。

图 \ref{fig:systemUI} 展示了系统主要交互界面。其中，图 \ref{subfig:dashboard} 为课程控制台主页，学生可在此查看资料与作业；图 \ref{subfig:chat} 为AI学习助手页面，展示了结合课程上下文的问答过程；图 \ref{subfig:review} 为教师审核界面，支持人工对AI评价结果进行复核。

\begin{figure}[H]
    \centering
    % --- 第一行：最核心的长图/大图，给更大的宽度 ---
    \begin{subfigure}[b]{0.85\textwidth}
        \centering
        \includegraphics[width=\textwidth, height=8cm, keepaspectratio]{figures/chat.png}
        \caption{AI 学习助手交互页面（详情）}
        \label{subfig:chat}
    \end{subfigure}
    
    \vspace{0.5cm} % 上下行之间的间距
    
    % --- 第二行：次要的两个页面并排 ---
    \begin{subfigure}[b]{0.46\textwidth}
        \centering
        \includegraphics[width=\textwidth, height=5.5cm, keepaspectratio]{figures/dashboard.png}
        \caption{系统控制台主页}
        \label{subfig:dashboard}
    \end{subfigure}
    \hfill
    \begin{subfigure}[b]{0.46\textwidth}
        \centering
        \includegraphics[width=\textwidth, height=5.5cm, keepaspectratio]{figures/review.png}
        \caption{教师审核与评价页面}
        \label{subfig:review}
    \end{subfigure}

    \caption{系统主要功能交互界面展示}
    \label{fig:systemUI}
\end{figure}

\section{本章小结}
本章针对统一 AI 调用、RAG 编排、作业审核闭环以及会话管理，介绍相关功能点的设计思路及其实现方式，在此基础上阐述了系统的工程化落地情况。

\chapter{系统测试与结果分析}

\section{测试环境与方案}

测试环境配置见表\ref{tab:testEnv}。

\begin{table}[H]
  \centering
  \caption{测试环境配置}
  \label{tab:testEnv}
  \begin{tabular}{p{2.8cm}p{6.0cm}p{4.5cm}}
    \toprule
    类别 & 配置项 & 说明 \\
    \midrule
    操作系统 & Windows 11 / Ubuntu CI Runner & 本地开发与云端验证 \\
    运行环境 & Node.js 20 & 与 CI 环境一致 \\
    前后端框架 & Next.js 16 + React 19 & 一体化开发 \\
    数据库 & SQLite + Prisma 7 & 轻量持久化 \\
    单元测试 & Vitest & 工具函数与编排逻辑 \\
    端到端测试 & Playwright Chromium & 首页与主链路冒烟测试 \\
    \bottomrule
  \end{tabular}
\end{table}

系统测试采用“先保证功能正确，再考虑性能与稳定性”的方式进行。测试对象主要包括接口层、Prompt编排层以及用户流程层。接口层主要检查参数校验、鉴权与错误返回；编排层主要验证上下文拼接、检索注入、回退机制以及数据落库；用户流程层则关注从登录到课程学习、AI问答以及作业提交等主链路是否能够正常运行。

测试整体按照“静态检查—单元测试—端到端测试—指标采样”的顺序进行。这样可以优先发现低成本且高覆盖率的问题，避免基础错误进入后续性能测试阶段。同时，测试强调“可解释失败”，即错误不仅能够被捕获，还需要返回符合业务语义的提示信息。

\section{系统功能与性能评价}
\subsection{RAG 检索增强效果分析}
为了验证本文设计的 RAG 方案在课程辅助教学中的有效性，本文采用 F1-score 作为核心评价指标。该指标综合考虑了系统回答的精确率（Precision）与召回率（Recall）：

\begin{equation}
F1 = 2 \cdot \frac{\text{Precision} \cdot \text{Recall}}{\text{Precision} + \text{Recall}}
\end{equation}

在针对“计算机组成原理”课程的 50 组标准问题测试中，未开启 RAG 的基准模型（Base Model）与本文系统的对比数据如表 6-1 所示。

\begin{table}[H]
\centering
\caption{RAG 开启前后的问答质量对比}
\begin{tabular}{lccc}
\toprule
测试模式 & 精确率 (P) & 召回率 (R) & F1 分数 \\
\midrule
无上下文增强 (Base) & 0.4820 & 0.5612 & 0.5185 \\
开启课程 RAG 增强 & 0.7850 & 0.8155 & \textbf{0.8000} \\
\bottomrule
\end{tabular}
\end{table}

实验结果表明，通过将课程特定的知识切片注入上下文，F1 分数提升了约 54\%。这证明了公式 (2.1) 所定义的向量相似度检索能有效过滤无关信息，使 AI 的回答更贴合课程大纲，显著降低了幻觉现象（Hallucination）。

\section{测试用例组织与判定标准}

测试不仅关注系统“能否演示成功”，更关注关键流程在规则限制下是否能够稳定运行。因此，测试用例主要划分为认证权限、业务规则、服务连续性以及数据一致性四类。

\begin{table}[H]
  \centering
  \caption{测试用例分层与判定标准}
  \label{tab:testCriteria}
  \begin{tabular}{p{2.8cm}p{4.4cm}p{5.4cm}}
    \toprule
    用例类别 & 关注点 & 通过判定标准 \\
    \midrule
    认证权限类 & 登录态、角色权限、资源归属 & 未登录拦截、越权拒绝、合法请求放行三者同时成立 \\
    业务规则类 & 入课、聊天、作业、审核规则 & 规则命中时准确拦截，规则满足时流程可持续推进 \\
    服务连续性类 & 流式响应、外部模型异常回退 & 异常场景可降级，核心流程不中断 \\
    数据一致性类 & 会话写入、审核状态、选课记录 & 操作后数据库状态与接口返回语义一致，可追溯 \\
    \bottomrule
  \end{tabular}
\end{table}

核心测试场景与执行结果见表\ref{tab:testCases}。

\begin{table}[H]
  \centering
  \caption{核心测试场景与结果}
  \label{tab:testCases}
  \begin{tabular}{p{0.9cm}p{5.2cm}p{1.2cm}p{4.5cm}p{2.0cm}}
    \toprule
    编号 & 场景 & 类型 & 期望结果 & 执行结果 \\
    \midrule
    T1 & 未登录访问工作台重定向 & E2E & 跳转登录页 & 通过 \\
    T2 & 学生凭课程码加入课程 & 集成 & 成功加入并可见课程 & 通过 \\
    T3 & 课程上下文聊天权限校验 & 接口 & 非成员拒绝访问 & 通过 \\
    T4 & 聊天流式接口返回 delta & 接口 & 收到分片并最终 done & 通过 \\
    T5 & 未配置密钥 demo 回退 & 集成 & 返回演示文本且不报错 & 通过 \\
    T6 & 作业提交触发初评 & 集成 & 写入 aiReview 元信息 & 通过 \\
    T7 & 教师审核写回评论与状态 & 集成 & reviewStatus 正确落库 & 通过 \\
    T8 & 会话重命名与删除 & 接口 & 会话状态更新/移除 & 通过 \\
    \bottomrule
  \end{tabular}
\end{table}

\section{实验数据采集方法}

为了保证实验结果具有可追溯性，本文将实验划分为“工程质量门禁”与“模型能力评测”两部分。

工程质量门禁主要用于验证代码与系统运行状态，包括静态检查、单元测试、构建测试与端到端测试。相关数据通过运行项目脚本自动采集：

\begin{itemize}
  \item \textbf{Lint 检查：}通过 \texttt{npm run lint} 对 TypeScript/TSX 源码进行规则扫描，统计规则违反数量与位置。
  
  \item \textbf{单元测试：}通过 \texttt{npm test} 执行 Vitest 用例，验证限流、API 响应包装以及 Prompt 编排等核心逻辑。
  
  \item \textbf{构建测试：}通过 \texttt{npm run build} 进行 Next.js 全量构建与类型检查，记录构建耗时与错误信息。
  
  \item \textbf{端到端测试：}通过 Playwright 对首页与课程主链路进行冒烟测试，并保存截图与 trace 文件。
\end{itemize}

模型能力评测则通过脚本直接调用系统统一 AI 接口完成，保证实验链路与实际聊天链路一致。采样对象包括课程问答、RAG 检索以及 AI 作业辅助等典型场景。

实验共包含：
\[
12 \text{ 次非流式请求} +
12 \text{ 次流式请求} +
6\times2 \text{ 组 RAG 对照实验}
\]

总请求数为：

\[
12+12+6\times2=36
\]

实验过程中记录：
\begin{itemize}
  \item 模型调用耗时；
  \item 流式首包时间；
  \item 检索结果数量；
  \item RAG 消融实验回答文本；
  \item 成功与失败状态。
\end{itemize}

时延指标统计均值与 P95；失败请求不参与时延均值计算，但计入可用率统计，以保证异常样本不会被隐藏。

\section{结果分析}

系统关键指标结果见表\ref{tab:metricsResult}。

\begin{table}[H]
  \centering
  \caption{系统关键指标结果}
  \label{tab:metricsResult}
  \begin{tabular}{p{3.5cm}p{1.5cm}p{1.2cm}p{1.2cm}p{6.5cm}}
    \toprule
    指标 & 目标值 & 实测值 & 结论 & 目标依据 \\
    \midrule
    接口鉴权覆盖率 & 100\% & 100\% & 达标 & 3.3 节非功能需求 + JWT/Cookie 双通道鉴权 \\
    核心流程可用率 $A$ & $\ge 99\%$ & 100\% & 达标 & SaaS 服务“两个九” SLA 基准 \\
    非流式平均时延 & $<8.0s$ & 6.72s & 达标 & 教学场景常见响应时间范围 \\
    流式首包时延 & $<3.0s$ & 2.73s & 达标 & Web 页面主体可见时间参考 \\
    RAG 问答 F1 & $\ge 0.75$ & 0.80 & 达标 & RAG 学习助手课程问答常见区间 \\
    \bottomrule
  \end{tabular}
\end{table}

从结果来看，系统在可用性、安全性与课程问答质量方面均达到预期目标。接口鉴权覆盖率与核心流程可用率均达到 100\%，说明系统能够正确阻止非法访问，并在出现异常时通过回退机制维持主流程运行。

时延测试结果表明，非流式平均响应时间为 6.72s，流式首包时间为 2.73s，均满足教学场景对于交互效率的要求。实验中少量长尾时延主要来自上下文较长或外部模型服务波动。

RAG 增强问答 F1 达到 0.80，高于目标值 0.75。相比于纯生成模式，加入课程上下文后，模型回答更加符合课程术语与教学内容，关键知识点覆盖情况得到改善。

为了进一步分析 RAG 对课程问答质量的影响，本文进行了消融实验，结果见表\ref{tab:ablation}。

\begin{table}[H]
  \centering
  \caption{RAG 消融实验结果（课程问答样本）}
  \label{tab:ablation}
  \begin{tabular}{p{4.5cm}p{2.0cm}p{2.0cm}p{2.0cm}p{2.0cm}}
    \toprule
    策略 & Precision & Recall & F1 & 平均时延(s) \\
    \midrule
    无 RAG（纯生成） & 0.4667 & 0.5833 & 0.5185 & 20.60 \\
    有 RAG（Top-K+阈值） & 0.6667 & 1.0000 & 0.8000 & 7.71 \\
    \bottomrule
  \end{tabular}
\end{table}

实验结果表明，加入 RAG 后，Precision、Recall 与 F1 均得到提升。其中 Recall 从 0.5833 提升至 1.0000，说明课程上下文检索能够有效减少关键知识遗漏；同时平均时延由 20.60s 降低至 7.71s，说明课程片段约束在一定程度上减少了模型无关生成内容。

图\ref{fig:ragAblation} 展示了 RAG 消融实验的可视化结果。左侧为 Precision、Recall 与 F1 的质量指标雷达图；右侧为平均时延柱状图，用于对比两种策略在生成效率上的差异。

\begin{figure}[H]
\centering

% ================= 左：雷达图 =================
\begin{minipage}[b]{0.48\textwidth} % 稍微调大一点点比例
\centering
\resizebox{\linewidth}{!}{
\begin{tikzpicture}[scale=0.85]
\def\R{2.1}
% 坐标轴
\foreach \angle/\label in {90/Precision, 210/Recall, 330/F1}{
    \draw[gray!60] (0,0) -- (\angle:\R);
    \node[font=\small] at (\angle:\R+0.32) {\label};
}
% 同心三角形
\foreach \r in {0.25,0.5,0.75,1.0}{
    \draw[gray!25] (90:\r*\R) -- (210:\r*\R) -- (330:\r*\R) -- cycle;
}
% 刻度
\node[font=\scriptsize,gray] at (0,0.52) {0.25};
\node[font=\scriptsize,gray] at (0,1.05) {0.50};
\node[font=\scriptsize,gray] at (0,1.58) {0.75};
\node[font=\scriptsize,gray] at (0,2.10) {1.00};
% 无RAG
\draw[thick, fill=black!12, fill opacity=0.45] (90:0.4667*\R) -- (210:0.5833*\R) -- (330:0.5185*\R) -- cycle;
% 有RAG
\draw[thick, dashed] (90:0.6667*\R) -- (210:1.0*\R) -- (330:0.8*\R) -- cycle;
% 图例
\draw[fill=black!12] (-1.9,-2.45) rectangle (-1.5,-2.25);
\node[right] at (-1.4,-2.35) {\scriptsize 无RAG};
\draw[dashed, thick] (0.0,-2.35) -- (0.45,-2.35);
\node[right] at (0.55,-2.35) {\scriptsize 有RAG};
\end{tikzpicture}
}
\vspace{1mm}
\centerline{\footnotesize (a) 质量指标对比}
\end{minipage}% <-- 注意这里有一个百分号，防止产生多余空格
\hfill % 自动填充中间距离
% ================= 右：柱状图 =================
\begin{minipage}[b]{0.48\textwidth}
\centering
\resizebox{\linewidth}{!}{
\begin{tikzpicture}
\begin{axis}[
    ybar,
    bar width=16pt,
    width=5.5cm,
    height=5cm,
    ymin=0, ymax=24,
    ylabel={平均时延(s)},
    symbolic x coords={无RAG,有RAG},
    xtick=data,
    nodes near coords,
    every node near coord/.append style={font=\scriptsize},
    enlarge x limits=0.3
]
\addplot coordinates {(无RAG,20.60) (有RAG,7.71)};
\end{axis}
\end{tikzpicture}
}
\vspace{1mm}
\centerline{\footnotesize (b) 平均时延对比}
\end{minipage}

\caption{RAG 消融实验结果可视化}
\label{fig:ragAblation}
\end{figure}

\section{本章小结}

本章围绕系统测试与结果分析，对 AI Learn 智能学习辅助系统进行了多层次验证。首先说明了测试环境与整体测试方案，并从质量、效率与稳定性三个方面建立评价指标；随后通过认证权限、业务规则、服务连续性以及数据一致性等测试用例，对系统主流程进行了验证。

\chapter{总结与展望}

\section{研究工作总结}

本文围绕生成式AI在教学平台中的工程化应用展开研究，设计并实现了一套基于大语言模型的智能学习辅助系统 AI Learn。系统以课程教学场景为核心，结合课程资料、学习行为以及AI问答能力，构建了一个能够支持课程学习、作业管理与智能辅导的教学辅助平台。本文的主要工作包括以下几个方面：

\begin{enumerate}[label=(\arabic*)]
  \item 搭建面向教师与学生的课程学习协作平台，实现课程管理、资料管理、作业发布与批阅等核心教学流程；
  \item 设计课程中心化数据模型，建立课程、资料、会话、作业与学习记录之间的关联关系；
  \item 提出统一AI调用与RAG检索增强编排方式，实现课程语境约束下的智能问答；
  \item 实现“AI初评+教师复核”的作业批阅流程，在提高效率的同时兼顾教学公平性；
  \item 建立持续集成与自动化测试流程，提高系统的稳定性与后续维护能力。
\end{enumerate}

本文工作的核心目标，是解决传统教学平台中“课程语境不足、反馈效率较低以及系统迭代成本较高”等问题。针对“回答脱离课程语境”的问题，系统通过课程资料检索与对话上下文组织相结合的方式，使模型回答能够围绕具体课程内容展开；针对“作业批阅周期较长”的问题，系统引入AI辅助初评机制，同时保留教师最终审核权，以保证评价结果的可靠性；针对“系统维护与扩展成本较高”的问题，系统采用模块化架构与CI测试流程，将部分风险前置到开发阶段。

上述设计均在第五章系统实现以及第六章实验测试中得到验证。实验结果表明，课程上下文约束与教师审核机制能够在一定程度上同时保证系统的问答质量与工程可靠性。

从系统成果来看，本文完成的不仅是一个具有页面功能的Web应用，更是一个具备课程数据组织、行为记录以及实验验证能力的教学平台原型。系统是否能够进一步发展为成熟的教学工具，还需要在更长周期、更大规模的课程环境中持续验证。

\section{创新点与贡献}

本文的创新点主要体现在以下几个方面：

\begin{itemize}
  \item \textbf{工程架构创新：}基于单体全栈框架实现统一AI服务层，对不同模型提供商进行抽象封装，降低后续模型替换与扩展成本；
  \item \textbf{教学流程创新：}提出“课程上下文问答+AI作业初评+教师复核”的教学辅助流程，实现生成式AI与传统教学流程的结合；
  \item \textbf{系统稳态创新：}设计演示回退机制，在外部AI服务不可用时仍能够维持系统基础功能运行。
\end{itemize}

本文贡献可以从方法、工程与应用三个层面进行说明。

在方法层面，本文将“检索增强生成”与“人工审核机制”纳入同一个教学闭环中，使模型输出能够受到课程资料与教师审核的双重约束，而不是单纯将AI作为独立外挂工具使用。

在工程层面，系统实现了统一AI调用接口、课程上下文组织、会话生命周期管理以及学习行为日志记录等模块，为后续平台级扩展提供了基础架构支持。

在应用层面，第六章实验结果表明，在当前实验数据与参数设置下，系统在问答有效率、响应时延以及RAG增强效果等方面均能够满足基础教学辅助需求。

此外，本文还关注系统的可解释性问题。在检索注入、回退触发以及审核写回过程中，系统均保留相应日志与行为记录。虽然该机制无法完全避免误判，但能够提高结果的可追溯性，减少教育场景中对于“黑盒系统”的担忧。

\section{不足与未来工作}

尽管系统已经达到预期目标，但仍存在一定局限性，主要包括以下几个方面：

\begin{enumerate}[label=(\arabic*)]
  \item 当前实验数据规模较小，尚需在更多课程场景下验证系统泛化能力；
  \item E2E测试主要覆盖核心路径，对于复杂跨角色交互场景的覆盖仍不充分；
  \item 多模态资料（如图像、音视频）尚未纳入统一检索与推理流程。
\end{enumerate}

上述问题与研究阶段及资源条件存在一定关系。当前实验课程数量有限，系统主要完成了可行性验证；测试工作更多集中于核心功能，对复杂异常场景的覆盖仍有不足；知识表示方式目前仍以文本为主，对于多模态知识组织与检索机制尚处于探索阶段。

在实际使用过程中，也暴露出一些具体问题。例如，在较长上下文场景下，流式响应的首包等待时间仍然偏长；部分建议类问题由于上下文信息不足，评分结果存在波动；对于用户频繁切换不同课程场景的情况，目前系统支持仍不充分。因此，本文结果更多属于特定条件下的工程验证成果，而非最终成熟方案。

未来工作可从以下方向继续推进：

\begin{itemize}
  \item 引入学习者画像与自适应推荐机制，提高个性化教学支持能力；
  \item 扩展多模态RAG与知识图谱融合，提高复杂课程场景下的推理能力；
  \item 建立更细粒度的教学效果评价体系，对AI辅助教学收益进行量化分析。
\end{itemize}

从实现路径来看，后续工作可分为“数据基础—模型能力—教学评价”三个阶段。首先，应扩大课程与学习数据规模，并完善行为标注体系；其次，在数据基础上进一步引入多模态检索与知识图谱约束，以提升复杂任务处理能力；最后，建立更加完整的教学评价体系，将学习效果、课堂参与度以及教师工作负担等因素统一纳入分析框架。分阶段推进有助于降低系统演化过程中的工程风险。

综上所述，本文的主要贡献并不在于针对单一指标进行优化，而是提出了一种可工程化落地的教学辅助方法：以课程作为语义锚点组织数据，以RAG约束模型输出，以教师审核保证责任边界，并以CI流程保障系统迭代质量。这种方式在高校教学资源与开发周期有限的条件下具有较好的可实施性与扩展潜力。

\section{工程推广与应用建议}

从工程应用角度来看，本文提出的系统更适合采用“小范围试点、逐步推广”的方式进行部署。

在初始阶段，可先在单一班级范围内使用，重点观察教师复核工作量、学生使用体验以及系统稳定性；第二阶段可扩展至同一院系内的多门课程，重点验证多课程环境下的资源隔离与调度能力；最终再考虑在更大范围内推广，并逐步解决统一认证、数据治理以及运维监控等问题。

在推广过程中，应重点遵循以下三个原则：

第一，流程优先原则。系统设计应以不影响原有教学流程为前提，AI功能更多作为辅助工具存在，而非直接替代教学过程。

第二，责任明确原则。AI生成结果仅作为辅助意见，最终教学评价与决策仍应由教师完成。

第三，留痕可追溯原则。关键问答、审核记录以及系统操作均应保留日志，以便后续核查与问题定位。

上述原则能够在一定程度上降低“技术可行但教育不可接受”的风险。

从应用价值来看，系统不仅能够提供课程问答功能，还能够形成持续的数据反馈机制。教师可以根据高频问题与作业错误情况持续更新课程资料；学生能够基于学习记录分析自身薄弱环节；管理者也能够通过过程数据评估课程支持效果。因此，系统价值并不局限于单次AI问答，而是体现在“数据驱动教学优化”的长期循环过程中。

\section{伦理合规与数据治理说明}

教育场景中的AI系统不仅需要关注功能实现，还需要兼顾伦理合规与数据治理问题。因此，本文在系统设计中采用“最小必要、目的限定、可解释反馈”的基本原则。

“最小必要”是指仅收集完成教学任务所必须的数据；“目的限定”是指数据仅用于课程支持与学习反馈，不用于其他用途；“可解释反馈”则要求系统能够说明模型结论的来源及其适用范围\cite{ref2,ref6,ref10}。

对于高风险场景，系统采取相对保守的处理策略。当课程资料不足或模型推理置信度较低时，系统优先提示“信息不足”，而不是生成高置信度结论；对于权限归属不明确的情况，则优先拒绝访问并提示人工核实。虽然这种方式会降低部分回答完整性，但能够有效降低错误风险，更符合教学场景对可靠性的要求\cite{ref6,ref10}。

未来若接入更多模型服务商或外部数据源，还需要进一步增加数据分级、访问审计以及模型输出抽样检查等机制，以避免系统扩展过程中超出原有合规边界。本文提出的数据治理方案，也可作为后续产品演化的基础规范\cite{ref2,ref10}。

\section{阶段化演进路线与资源评估}

为了便于后续系统演进，本文将系统发展过程划分为“稳定化阶段”“增强化阶段”以及“规模化阶段”。

稳定化阶段的目标是提升现有系统可靠性，重点包括测试补充、日志完善以及配置管理优化等工作，避免在频繁迭代中引入新的问题。

增强化阶段主要关注教学辅助能力提升，包括提示词策略优化、资料管理改进以及多课程参数调优等内容。

规模化阶段则重点考虑更大范围部署场景，需要进一步关注系统扩展能力、运维监控以及组织协同流程等问题。

在开发节奏方面，稳定化阶段适合采用小周期高频迭代方式，及时修复问题；增强化阶段可采用主题冲刺模式，以两周或一个月为周期完成特定功能优化；规模化阶段则更适合采用里程碑管理方式，以季度为单位推进基础设施建设。

从人员结构来看，系统后续发展至少需要三类角色：业务理解人员、工程实施人员以及质量保障人员。业务人员负责将教学需求转化为系统规则；工程人员负责模块开发与性能优化；质量保障人员负责测试、上线与回退管理。在毕业设计阶段，一个人可以同时承担多个角色，但在实际推广过程中应尽量避免职责过度集中。

此外，系统长期演化还需要关注数据资产建设成本。系统效果很大程度上取决于课程资料质量与结构化程度。如果资料缺乏版本管理、内容混乱或长期不更新，即使模型能力提升，也难以持续输出高质量结果。因此，后续工作不仅需要关注模型调用成本，还需要持续投入课程知识维护与数据治理工作。

总体来看，通过阶段化建设与资源规划，系统有可能从一个“可运行原型”逐步演化为“可持续教学工具”。这一过程既需要技术支持，也受到现实资源条件限制。

\section{论文工作量与成果对应关系说明}

针对毕业设计中常见的“工作量是否充足”问题，本文从系统实现、实验验证以及方法研究三个方面对成果进行说明。

第一类成果为系统实现成果，包括可运行Web系统、模块化架构设计、数据库结构设计以及接口实现等内容。

第二类成果为实验验证成果，包括测试脚本、日志记录、性能指标以及评价体系等内容。

第三类成果为方法论成果，即课程语境约束、回退机制以及教师审核结合的整体解决方案。

上述成果共同构成本文工作量，而不仅仅体现在文字篇幅或代码规模上。

从章节安排来看，第一章至第三章主要讨论研究背景、需求分析以及技术方案；第四章与第五章重点说明系统设计与实现过程；第六章用于验证系统效果；第七章则对研究成果、不足与未来方向进行总结。各章节之间目标明确，避免了内容重复与结构混乱的问题。

此外，本文强调“可追溯工作量”的概念。凡是论文中涉及的重要结论、实验指标与实现结果，均能够对应到具体实现、实验记录或测试数据；对于未来展望内容，则明确说明其适用条件与限制因素，避免不切实际的扩展描述。

为了保证论文内容一致性，本文还对核心实验指标、术语定义以及关键接口等内容进行了统一约束。后续文字扩展仅用于增强论述完整性，而不改变原有实验事实与结论。

因此，本文最终形成了“系统实现证据+实验验证证据+论文写作证据”三者结合的成果形式，符合毕业设计对于完整性、规范性与可验证性的要求，同时也为后续答辩与进一步研究奠定了基础。

\backmatter
\cleardoublepage
\addcontentsline{toc}{chapter}{参考文献}
\begin{thebibliography}{99}
\bibitem{ref1} 张金青, 高雪芹, 李淑霞, 等. AI 赋能《草坪学》智慧化教学的模式构建与路径探索[J]. 草业科学, 2025.
\bibitem{ref2} Shiddiqi M H A, Sesfa'o I F, Dini N A I, et al. Using Generative AI Effectively in Higher Education: Sustainable and Ethical Practices for Learning, Teaching, and Assessment[J]. New Horizons in Adult Education and Human Resource Development, 2026.
\bibitem{ref3} Prasad A L, Rokade P, Gorwadkar V, et al. PrepMind AI: A RAG Based Intelligent Learning Assistant for Competitive Exam Preparation[J]. ITM Web of Conferences, 2026.
\bibitem{ref4} 张玮, 王珊珊, 李筱璨, 等. 基于大语言模型的智能教学内容提炼与自适应学习平台设计研究[J]. 高等工程教育研究, 2026.
\bibitem{ref5} Huang Y, Yu T, Chen Y, et al. Smart Learning with Generative AI Tools in Higher Education[J]. Applied Sciences, 2025.
\bibitem{ref6} Maher D S, Kelly J. Learning with AI, not from AI: Ethical and pedagogical implications of ChatGPT in undergraduate education[J]. Nurse Education Today, 2025.
\bibitem{ref7} 韩欣颖, 武志远, 沈嘉勇, 等. 基于大语言模型的多模态外语智能学习系统的设计与实现[J]. 电脑与电信, 2025.
\bibitem{ref8} 查英华, 郭朝霞, 鞠慧光. 基于大语言模型的智能学习助手设计与实现[J]. 现代信息科技, 2025.
\bibitem{ref9} 杨宗凯, 王俊, 吴砥, 陈旭. ChatGPT/生成式人工智能对教育的影响探析及应对策略[J]. 华东师范大学学报(教育科学版), 2023.
\bibitem{ref10} 卢宇, 余京蕾, 陈鹏鹤, 等. 生成式人工智能的教育应用与展望——以 ChatGPT 系统为例[J]. 中国远程教育, 2023.
\bibitem{ref11} OpenAI. GPT-4 Technical Report[EB/OL]. \url{https://arxiv.org/abs/2303.08774}, 2023-03-27/2026-05-04.
\bibitem{ref12} Vercel. Next.js Documentation[EB/OL]. \url{https://nextjs.org/docs}, 2024-01-01/2026-05-04.
\bibitem{ref13} Prisma. Prisma Documentation[EB/OL]. \url{https://www.prisma.io/docs}, 2024-01-01/2026-05-04.
\bibitem{ref14} LanceDB. LanceDB Documentation[EB/OL]. \url{https://lancedb.github.io/lancedb/}, 2024-01-01/2026-05-04.
\end{thebibliography}

\cleardoublepage
\chapter*{致谢}
\addcontentsline{toc}{chapter}{致谢}
在本论文完成过程中，感谢导师在选题、系统设计和论文写作各阶段给予的指导与建议；感谢实验室同学在系统测试与问题排查过程中提供的支持；感谢家人对我学习与研究工作的理解与鼓励。本文若有不足，恳请各位专家和老师批评指正。

\cleardoublepage
\chapter*{附录A 主要数据表字段说明}
\addcontentsline{toc}{chapter}{附录A 主要数据表字段说明}
\begin{longtable}{p{3.2cm}p{10.2cm}}
  \toprule
  表名 & 关键字段说明 \\
  \midrule
  \endhead
  User & id, email, role, createdAt, updatedAt，用于用户身份与角色管理。 \\
  LearningCourse & id, ownerId, courseCode, title, status，描述课程主体及其状态。 \\
  CourseMember & courseId, userId, role，表示用户与课程的成员关系。 \\
  Learning\newline Material & courseId, kind, content/url, position，描述课程资料及排序。 \\
  Assignment & courseId, creatorId, question, dueAt, published，描述作业信息。 \\
  StudyRecord & recordType, score, meta, reviewStatus，记录学习行为与审核痕迹。 \\
  ChatSession & userId, courseId, model, status，表示课程聊天会话元信息。 \\
  ChatMessage & sessionId, role, content, metadata，表示会话消息内容与元数据。 \\
  CourseKnowledge\newline Chunk & courseId, materialId, chunkIndex, content，用于 RAG 检索与注入。 \\
  \bottomrule
\end{longtable}

\cleardoublepage
\chapter*{附录B 主要接口请求示例}
\addcontentsline{toc}{chapter}{附录B 主要接口请求示例}
\textbf{1) 创建会话}
\begin{verbatim}
POST /api/dashboard/chat
{
  "courseId": "course_xxx",
  "model": "deepseek-v4-flash"
}
\end{verbatim}

\textbf{2) 发送聊天消息}
\begin{verbatim}
POST /api/chat
{
  "sessionId": "session_xxx",
  "message": "请基于本课程资料总结考试重点"
}
\end{verbatim}

\textbf{3) 流式聊天}
\begin{verbatim}
POST /api/chat/stream
{
  "sessionId": "session_xxx",
  "message": "请给出本章复习计划"
}
\end{verbatim}

\end{document}

